\section{3D生成方法的演进}

\subsection{早期方法：3D GAN（2016）}

最早的生成式3D建模方法使用生成对抗网络（GAN）。在2016年的工作中，研究者构建了一个纯生成模型：从噪声隐变量出发，训练模型生成3D体素（voxel）表示，并使用\textbf{3D判别器}进行对抗训练。

该方法的两个主要限制：
\begin{itemize}
    \item \textbf{无法精确控制输入}：生成过程完全随机，无法指定生成特定物体
    \item \textbf{需要3D真值数据}：训练依赖于3D ground truth，而高质量3D数据极为稀缺
\end{itemize}

\subsection{PlatonicGAN：摆脱3D数据依赖}

PlatonicGAN 的核心创新在于引入了\textbf{可微渲染器}（differentiable renderer）和\textbf{2D判别器}，从而消除了对3D真值数据的依赖。这使得模型可以直接在互联网上的2D图像数据上进行训练。此外，PlatonicGAN 还引入了\textbf{图像编码器}，实现了从单张图像进行3D重建的能力。

\begin{knowledgebox}{可微渲染的意义}
可微渲染（differentiable rendering）是连接3D表示与2D监督信号的桥梁。通过使渲染过程可微分，梯度可以从2D图像空间的损失函数反向传播到3D表示的参数上，从而允许使用大量的2D图像数据来训练3D生成模型，而无需昂贵的3D标注。
\end{knowledgebox}

\subsection{PixelNeRF：多视角融合}

在基于图像的重建思路基础上，PixelNeRF 专注于\textbf{多视角输入}。其核心创新是将来自多张输入图像的编码特征聚合（aggregate），融合为一个一致的3D表示——具体形式为 NeRF（Neural Radiance Field）。

\subsection{共同特征与局限}

上述所有方法都有一个共同特点：它们直接预测一个\textbf{显式的、3D一致的}表示。

\begin{importantbox}{显式3D表示的优势与劣势}
\textbf{优势}：一旦获得3D模型，可以从任意视角快速、低成本地渲染，非常适合手机和平板等移动端的实时应用。

\textbf{劣势}：整体生成质量较低。原因有二：（1）高质量3D数据极度稀缺；（2）这些方法不应使用过强的归纳偏置（如渲染），但避免归纳偏置又限制了模型能力。
\end{importantbox}

\subsection{本章小结}

3D生成方法从依赖3D真值数据的3D GAN，发展到利用可微渲染和2D判别器的 PlatonicGAN，再到支持多视角输入的 PixelNeRF。这些方法的共同优势在于生成显式3D表示、渲染成本低，但受限于3D数据的稀缺性，整体质量不及2D生成模型。

